\section{Beyond log-concavity}
\label{sec:nonconvex}

Strong log-concavity plays two roles in the preceding theory: it supplies a
global functional inequality and it keeps the covariance equation away from
negative-curvature amplification.  A globally bounded Hessian alone supplies
neither.  We make the failure explicit on a smooth one-dimensional target and
then recover a constrained stationarity statement by clipping the covariance.

\subsection{A smooth double well}
\label{sec:nonconvex-double-well}

For \(R>0\), consider
\begin{equation}
  V_R(x)
  :=
  \frac{x^2}{2}-2R^2\log\cosh(x/R).
  \label{eq:nonconvex-potential}
\end{equation}
Its derivatives are
\[
  V_R'(x)=x-2R\tanh(x/R),
  \qquad
  V_R''(x)=1-2\operatorname{sech}^2(x/R),
\]
so \(-1\leq V_R''\leq1\) uniformly in \(R\).  For
\(X\sim\N(0,c)\), write
\[
  A_R(c):=\E[V_R''(X)].
\]
For every fixed \(c>0\), dominated convergence gives
\begin{equation}
  A_R(c)\longrightarrow-1
  \quad\text{as }R\to\infty.
  \label{eq:nonconvex-negative-curvature-limit}
\end{equation}
The symmetric mean \(m=0\) is invariant for all the covariance dynamics below.

\begin{proposition}[Riemannian covariance cascade]
\label{prop:nonconvex-riemannian-cascade}
Fix \(h>0\), and let
\begin{equation}
  c_{n+1}^{(R)}
  =
  c_n^{(R)}
  \exp\!\left\{h\bigl(1-c_n^{(R)}A_R(c_n^{(R)})\bigr)\right\},
  \qquad c_0^{(R)}=1.
  \label{eq:nonconvex-riemannian-map}
\end{equation}
For every \(M>0\), there are \(R\) and a finite \(n\) such that
\(c_n^{(R)}\geq M\).  Consequently, no bound depending only on
\(\sup_x|V_R''(x)|\), \(h\), and \(c_0\) can control all trajectories of the
unprojected Riemannian scheme uniformly over this family.
\end{proposition}

\begin{proof}
The limiting recursion obtained from
\eqref{eq:nonconvex-negative-curvature-limit} is
\[
  \bar c_{n+1}=\bar c_n e^{h(1+\bar c_n)},\qquad \bar c_0=1,
\]
and diverges.  Choose a finite \(N\) with \(\bar c_N>2M\).
The map \((R,c)\mapsto A_R(c)\) converges uniformly to \(-1\) on every
compact \(c\)-interval.  An induction over the first \(N\) steps therefore
gives \(c_n^{(R)}\to\bar c_n\) for \(0\leq n\leq N\).
For sufficiently large \(R\), \(c_N^{(R)}>M\).
\end{proof}

\begin{proposition}[Pole of the KL covariance resolvent]
\label{prop:nonconvex-kl-pole}
For \(h=1\), the unprojected KL covariance map is
\begin{equation}
  c_1
  =
  \frac{2c_0}{1+c_0A_R(c_0)}.
  \label{eq:nonconvex-kl-map}
\end{equation}
There are sequences \(\varepsilon\downarrow0\) and \(R_\varepsilon\to\infty\)
such that, with \(c_0=1-\varepsilon\),
\[
  c_1=\Theta(\varepsilon^{-1}),
  \qquad
  \bigl(1-c_1A_{R_\varepsilon}(c_1)\bigr)^2
  =\Theta(\varepsilon^{-2}).
  \label{eq:nonconvex-kl-pole-scaling}
\]
Thus the pole is caused by approaching loss of positive definiteness in the
rational resolvent, despite the uniform Hessian bound.
\end{proposition}

\begin{proof}
For every \(c>0\), writing \(X=\sqrt c Z\) gives the quantitative bound
\begin{equation}
  0\leq A_R(c)+1
  =2\E\tanh^2(\sqrt c Z/R)
  \leq \frac{2c}{R^2}.
  \label{eq:nonconvex-curvature-quantitative}
\end{equation}
Take \(0<\varepsilon\leq1/4\) and choose, for example,
\(R_\varepsilon^2\geq2\varepsilon^{-2}\).  Then
\(0\leq A_{R_\varepsilon}(1-\varepsilon)+1\leq\varepsilon^2\), and the
denominator in
\eqref{eq:nonconvex-kl-map} is then
\[
  1-(1-\varepsilon)+O(\varepsilon^2)
  =\varepsilon+O(\varepsilon^2),
\]
which is positive and gives \(c_1=\Theta(\varepsilon^{-1})\); in particular,
\(c_1\leq2/\varepsilon\).  Applying
\eqref{eq:nonconvex-curvature-quantitative} once more,
\[
  0\leq A_{R_\varepsilon}(c_1)+1
  \leq \frac{2c_1}{R_\varepsilon^2}
  \leq2\varepsilon\leq\frac12.
\]
This curvature bound yields
\(\abs{1-c_1A_{R_\varepsilon}(c_1)}=\Theta(c_1)\), proving the
residual scaling.
\end{proof}

These propositions concern the two discrete maps.  They do not say that every
bounded-Hessian nonconvex target diverges, nor do they preclude local stability
around a favorable stationary point.  They show that a global smoothness
constant cannot replace covariance control in a trajectory-wise
Fisher--Rao theorem.

\subsection{A Bures--Wasserstein comparison}
\label{sec:nonconvex-bw}

The Bures--Wasserstein forward--backward method has a different nonconvex
guarantee.  Under the assumptions and stepsize of
\citet[Theorem~5.5]{diao2023forward}, its iterates satisfy a running-minimum
bound of the form
\begin{equation}
  \min_{0\leq n<N}
  \norm{\grad_{\BW}\calE(a_n)}_{\BW}^2
  \leq
  \frac{150\{\calE(a_0)-\inf\calE\}}{\eta N}.
  \label{eq:nonconvex-bw-running-min}
\end{equation}
This is a stationarity statement in the Bures--Wasserstein geometry.  It does
not control the pointwise Fisher--Rao residual from
\Cref{prop:nonconvex-riemannian-cascade,prop:nonconvex-kl-pole}.
We use it only as a comparison between geometries, not as the first stage of a
general WFR method.

\subsection{Covariance-constrained KL stationarity}
\label{sec:nonconvex-clipped}

Fix \(0<\lambda_-\leq\lambda_+<\infty\) and
\[
  \mathcal A_{[\lambda_-,\lambda_+]}
  :=
  \{(m,C):\lambda_-\Id\preceq C\preceq\lambda_+\Id\}.
\]
For a symmetric matrix \(C=Q\diag(c_i)Q^\top\), let
\[
  \operatorname{Clip}_{[\lambda_-,\lambda_+]}(C)
  :=
  Q\diag\!\left(
    \min\{\lambda_+,\max\{\lambda_-,c_i\}\}
  \right)Q^\top.
\]
Assume only
\begin{equation}
  -\beta\Id\preceq\nabla^2V(x)\preceq\beta\Id
  \qquad\text{for all }x.
  \label{eq:nonconvex-smoothness}
\end{equation}
At \(a_n=(m_n,C_n)\), set
\[
  b_n:=\E_{\N(m_n,C_n)}\nabla V(X),
  \qquad
  A_n:=\E_{\N(m_n,C_n)}\nabla^2V(X),
\]
and apply the split update
\begin{align}
  m_{n+1}
  &=m_n-hC_nb_n,
  \label{eq:nonconvex-clipped-mean}\\
  \widetilde C_{n+1}
  &=(1+h)(C_n^{-1}+hA_n)^{-1},
  \qquad
  C_{n+1}
  =\operatorname{Clip}_{[\lambda_-,\lambda_+]}
    (\widetilde C_{n+1}).
  \label{eq:nonconvex-clipped-covariance}
\end{align}
The inverse is well defined at the stepsizes below.  The natural one-step
stationarity measure for this split scheme is
\begin{equation}
  \mathsf S_n
  :=
  \frac12\norm{m_{n+1}-m_n}_{C_n^{-1}}^2
  +
  \KL\!\left(
    \N(0,C_n)\,\middle\|\,\N(0,C_{n+1})
  \right).
  \label{eq:nonconvex-split-stationarity}
\end{equation}
The covariance orientation in the KL term is important.

\begin{theorem}[Constrained split stationarity]
\label{thm:nonconvex-clipped-stationarity}
Suppose \eqref{eq:nonconvex-smoothness} holds,
\(a_0\in\mathcal A_{[\lambda_-,\lambda_+]}\), and
\[
  0<h\leq\frac1{2\beta\lambda_+}.
\]
Then the iterates
\eqref{eq:nonconvex-clipped-mean}--\eqref{eq:nonconvex-clipped-covariance}
are well defined, remain feasible, and satisfy
\begin{equation}
  \calE(a_{n+1})
  \leq
  \calE(a_n)-\frac1h\mathsf S_n.
  \label{eq:nonconvex-clipped-descent}
\end{equation}
Consequently, for every \(N\geq1\),
\begin{equation}
  \boxed{
  \min_{0\leq n<N}\mathsf S_n
  \leq
  \frac{h}{N}
  \bigl\{\calE(a_0)-\calE(a_N)\bigr\}.}
  \label{eq:nonconvex-clipped-running-min}
\end{equation}
\end{theorem}

\begin{proof}
The spectral bounds imply
\[
  C_n^{-1}+hA_n
  \succeq
  (\lambda_+^{-1}-h\beta)\Id
  \succeq(2\lambda_+)^{-1}\Id.
\]
The smooth part of the Gaussian objective is relatively smooth with respect
to the forward split divergence with constant \(2\beta\lambda_+\).  The
mean optimality equation leaves
\(\frac12\norm{m_{n+1}-m_n}_{C_n^{-1}}^2\), while the constrained
log-determinant three-point inequality leaves the reverse covariance
divergence
\(\KL(\N(0,C_n)\|\N(0,C_{n+1}))\).  Spectral clipping is precisely the
log-determinant Bregman projection onto the covariance interval.  Since
\(h(2\beta\lambda_+)\leq1\), the forward remainder is absorbed and the two
estimates combine to
\eqref{eq:nonconvex-clipped-descent}.  Summing over \(n\) and taking the
smallest term gives \eqref{eq:nonconvex-clipped-running-min}.
Full details are recorded in \Cref{app:nonconvex-details}.
\end{proof}

\begin{corollary}[Full successive-Gaussian KL]
\label{cor:nonconvex-full-kl}
Under the assumptions of
\Cref{thm:nonconvex-clipped-stationarity},
\begin{equation}
  \min_{0\leq n<N}
  \KL\!\left(
    \N(m_n,C_n)\,\middle\|\,\N(m_{n+1},C_{n+1})
  \right)
  \leq
  \frac{\lambda_+}{\lambda_-}
  \frac{h}{N}
  \{\calE(a_0)-\calE(a_N)\}.
  \label{eq:nonconvex-full-kl-bound}
\end{equation}
For centered scalar trajectories the mean term vanishes, so the factor-one
bound \eqref{eq:nonconvex-clipped-running-min} is the quantity plotted in
\Cref{sec:numerics-nonconvex}.
\end{corollary}

\begin{proof}
The reverse KL between successive Gaussians splits into the covariance KL in
\eqref{eq:nonconvex-split-stationarity} and
\(\frac12\norm{m_{n+1}-m_n}_{C_{n+1}^{-1}}^2\).
Feasibility gives
\[
  C_{n+1}^{-1}
  \preceq\lambda_-^{-1}\Id
  \preceq\frac{\lambda_+}{\lambda_-}C_n^{-1}.
\]
Therefore the full KL is at most
\((\lambda_+/\lambda_-)\mathsf S_n\), and the conclusion follows from
\Cref{thm:nonconvex-clipped-stationarity}.
\end{proof}

\begin{remark}[Scope]
\label{rem:nonconvex-scope}
The result certifies stationarity of the projected split map on the
covariance-truncated set.  If an iterate is pinned at
\(\lambda_+\Id\), \(\mathsf S_n\) can vanish while the unconstrained
Fisher--Rao gradient points outward.  The theorem therefore does not imply
small unconstrained gradient, convergence to the unconstrained Gaussian
optimizer, or a stochastic clipped guarantee.
\end{remark}
